%%%%%%%%%%%%%%%%%%%%%%%%%%%%%%%%%%%%%%%%%%%%%%%%%%%%%%%%%%%%%%%%%%%%%%%%%%%%%%%%

\cleardoubleoddpage
\begin{abstract*}
------------------------ \textbf{PROTOTYPE - NOT FINAL} ------------------------  \break
As the demand for higher refresh-rates in video-games rises due to the increasing adoption of high-refresh-rate display panels, developers are restricted by the plateauing single-core performance and memory-latency of user-grade CPUs, as well as higher expectations towards the final product in areas such as more complex gameplay mechanics and large-scale physics simulations. 

To combat this widening gap, more and more game-engines and development studios employ Data Oriented Programming (DOP) approaches instead of classical Object Oriented Programming (OOP) approaches. This is due to DOP using less abstraction than OOP and, therefore, creating code that is easier to compute and optimize, for example, by using single instruction, multiple data (SIMD) vectorization. This thesis aims to quantify the performance difference between the two approaches using the Bevy Engine in Rust. Multiple DOP and  reference OOP benchmarks are implemented and described, investigating scenarios in which DOP outperforms OOP and where the classical OOP approach may prove more suitable. 
\end{abstract*}

%%%%%%%%%%%%%%%%%%%%%%%%%%%%%%%%%%%%%%%%%%%%%%%%%%%%%%%%%%%%%%%%%%%%%%%%%%%%%%%%

\cleardoubleoddpage
\begin{otherlanguage}{ngerman}
\begin{abstract*}
< Abstract auf Deutsch >
\end{abstract*}
\end{otherlanguage}

%%%%%%%%%%%%%%%%%%%%%%%%%%%%%%%%%%%%%%%%%%%%%%%%%%%%%%%%%%%%%%%%%%%%%%%%%%%%%%%%
